Can an ant on a surface tell that it is curved? Why does a full turn of an electron not bring it back to where it started? And
why do the eigenvalues of a random matrix space themselves exactly like the zeros of the Riemann zeta function?

This book answers with the geometry of curved spaces. Manifolds, flows and differential forms, with Stokes' theorem and Brouwer's
fixed point theorem in every dimension; curvature, from Gauss's Theorema Egregium to the Gauss--Bonnet theorem; the fundamental
group, covering spaces and homology, which prove that spaces of different dimensions are never the same. Then the shapes that
groups make: matrix Lie groups and the double cover of the rotations, Haar measure and Weyl's law for the eigenvalues of random
unitary matrices, the hyperbolic plane where continued fractions are paths through a tiling, and integral geometry, where
Buffon's needle measures the length of every curve.

Every theorem is proved in full, every exercise has a hint and a reasoned solution, and every chapter comes with a Python
notebook that puts the ideas to work. Written for the pleasure of seeing why things are true.
